\documentclass[border=4pt]{standalone}
\usepackage{amsmath,amssymb,mathtools,xcolor,varwidth}
\definecolor{subred}{HTML}{D81B60}
\definecolor{arcgreen}{HTML}{00A35A}
\definecolor{tanblue}{HTML}{1E6FE0}
\definecolor{auxpurple}{HTML}{8E24AA}
\begin{document}
\begin{varwidth}{28em}
\large\bfseries
Drop \textcolor{subred}{$BD$ perpendicular to the tangent} — this is the "subtense of the      angle of contact," the precise measure of how far the \textcolor{arcgreen}{curve} has risen      above \textcolor{tanblue}{$AD$}. Now the clever auxiliary: draw $BG\perp AB$ and $AG\perp AD$,      meeting at $G$. Because both $\angle ABG$ and $\angle AGD$ are right, the points      $A,B,G$ lie on \textcolor{auxpurple}{a circle with $AG$ as its diameter}. A chord dropped      from that circle gives the exact relation $AB^{2}=AG\times BD$ — trapping the tiny      \textcolor{subred}{departure $BD$} inside a clean, computable ratio.
\end{varwidth}
\end{document}
